%!TEX TS-program = xelatex
\documentclass{nihbiosketch}

%------------------------------------------------------------------------------
\name{Barrows, Anthony J.}
\eracommons{abarrows}
\position{PhD Student in Complex Systems and Data Science}

\begin{document}
	%------------------------------------------------------------------------------
	\begin{education}
		University of Hartford, Hartford, CT & BA & 05/2016 & Psychology \\
		University of Vermont, Burlington, VT & MS & 05/2021 & Complex Systems and Data Science \\
		University of Vermont, Burlington, VT & PhD & Expected 2026 & Complex Systems and Data Science \\
	\end{education}
	
	\section{Personal Statement}
	
	My research focuses on understanding brain-behavior relationships through computational and predictive modeling, with particular emphasis on cognitive control processes and substance use behaviors. I am currently pursuing a PhD in Complex Systems and Data Science at the University of Vermont, where I have developed expertise in applying cognitive process models to reveal meaningful associations between neural activation patterns and behavioral outcomes. My work primarily uses data from large-scale developmental studies, including the Adolescent Brain Cognitive Development (ABCD) Study, to investigate how cognitive mechanisms underlying behavioral tasks relate to brain function.
	
	My research interests span the intersection of computational neuroscience, cognitive psychology, and public health. I am particularly interested in how individual differences in brain function and cognition relate to risk factors for psychopathology and substance use. Through my position as a NERVE Lab predoctoral fellow, I have gained extensive experience in neuroimaging data analysis, statistical learning, and computational modeling techniques, as well as a deep commitment to open science and open-source software. 
	
	My long-term career goals center around using robust, reproducible analytical techniques to understand the neurocognitive mechanisms underlying behavioral health outcomes, with the ultimate aim of informing intervention strategies. I am committed to advancing methodological best practices and promoting the adoption of reproducible research workflows that can be easily shared, validated, and extended by the broader scientific community.
	
	Citations: 
	
	\begin{enumerate}
		\item Barrows, A., Weigard, A., McCabe, M., Potter, A., Garavan, H., and Allgaier, N. (2025). Cognitive process models reveal meaningful brain-behavior associations in the abcd study stop-signal task. \emph{Submitted to Journal of Cognitive Neuroscience}
		\item Barrows, A., Klemperer, E., Garavan, H., Allgaier, N., Lindson, N., and Taylor, G. (2025). Smoking reduction trajectories, and their association with smoking cessation: A secondary analysis of longitudinal RCT data. \emph{Submitted to BMJ Public Health}
		\item Higgins, S. T., Tidey, J. W., Sigmon, S. C., Heil, S. H., Gaalema, D. E., Lee, D., Hughes, J. R., Villanti, A. C., Bunn, J. Y., Davis, D. R., Bergeria, C. L., Streck, J. M., Parker, M. A., Miller, M. E., DeSarno, M., Priest, J. S., Cioe, P., MacLeod, D., Barrows, A. J., Markesich, C., and Harfmann, R. F. (2020). Changes in Cigarette Consumption With Reduced Nicotine Content Cigarettes Among Smokers With Psychiatric Conditions or Socioeconomic Disadvantage: 3 Randomized Clinical Trials. \emph{JAMA Network Open}, 3(10):e2019311--e2019311
	\end{enumerate}
	
	%------------------------------------------------------------------------------
	\section{Positions and Honors}
	\subsection*{Positions and Employment}
	\begin{datetbl}
		2022--Present & NERVE Lab Predoctoral Fellow (T32 DA043593) \\
		2022--Present & PhD Student, Complex Systems and Data Science Program, University of Vermont \\
		2026--2022 & Research Assistant and Data Analyst, Vermont Center on Behavior and Health \\
	\end{datetbl}
	
	%------------------------------------------------------------------------------
	\section{Contribution to Science}
	
	\begin{enumerate}
		\item \textbf{Computational Modeling of Cognitive Control Processes}
		
		My current work involves applying computational models to understand brain-behavior relationships in cognitive control tasks. I have focused on the stop-signal task as a paradigm for measuring response inhibition, utilizing cognitive process models which decompose behavioral performance into underlying cognitive mechanisms. This approach has revealed more meaningful associations between neural activation patterns and behavioral outcomes compared to traditional behavioral measures.
		
		\begin{enumerate}
			\item Barrows, A., Weigard, A., McCabe, M., Potter, A., Garavan, H., and Allgaier, N. (2025). Cognitive process models reveal meaningful brain-behavior associations in the abcd study stop-signal task. \emph{Submitted to Journal of Cognitive Neuroscience}
		\end{enumerate}
		
		\item \textbf{Longitudinal Analysis of Smoking Behavior and Cessation}
		
		I have contributed to understanding smoking reduction patterns and their relationship to cessation outcomes through secondary analysis of longitudinal randomized controlled trial data. This work has implications for developing more effective smoking cessation interventions by quantifying the initial reduction patterns most predictive of later cessation.
		
		\begin{enumerate}
			\item Barrows, A., Klemperer, E., Garavan, H., Allgaier, N., Lindson, N., and Taylor, G. (2023). Smoking reduction trajectories, and their association with smoking cessation: A secondary analysis of longitudinal RCT data. \emph{Submitted to BMJ Public Health}
		\end{enumerate}
		
		\item \textbf{Reduced Nicotine Content Cigarettes Research}
		
		I have contributed to multiple studies examining the effects of reduced nicotine content cigarettes on smoking behavior among various populations, including those with psychiatric conditions and socioeconomic disadvantage. This research contributes to tobacco regulatory science and potential policy interventions.
		
		\begin{enumerate}
			\item Higgins, S. T., Tidey, J. W., Sigmon, S. C., Heil, S. H., Gaalema, D. E., Lee, D., Hughes, J. R., Villanti, A. C., Bunn, J. Y., Davis, D. R., Bergeria, C. L., Streck, J. M., Parker, M. A., Miller, M. E., DeSarno, M., Priest, J. S., Cioe, P., MacLeod, D., Barrows, A. J., Markesich, C., and Harfmann, R. F. (2020). Changes in Cigarette Consumption With Reduced Nicotine Content Cigarettes Among Smokers With Psychiatric Conditions or Socioeconomic Disadvantage: 3 Randomized Clinical Trials. \emph{JAMA Network Open}, 3(10):e2019311--e2019311
			\item Nighbor, T. D., Barrows, A. J., Bunn, J. Y., DeSarno, M. J., Oliver, A. C., Coleman, S. R. M., Davis, D. R., Streck, J. M., Reed, E. N., Reed, D. D., and Higgins, S. T. (2020). Comparing Participant Estimated Demand Intensity on the Cigarette Purchase Task to Consumption when Usual-Brand Cigarettes were Provided Free. \emph{Preventive Medicine}
			\item Parker, M. A., Streck, J. M., Bergeria, C. L., Bunn, J. Y., Gaalema, D. E., Davis, D. R., Barrows, A. J., Sigmon, S. C., Tidey, J. W., Heil, S. H., and Higgins, S. T. (2018). Reduced Nicotine Content Cigarettes and Cannabis Use in Vulnerable Populations. \emph{Tobacco Regulatory Science}, 4(5):94--91
		\end{enumerate}
	\end{enumerate}
	
	\subsection*{Complete List of Published Work in MyBibliography:}
	\url{https://www.ncbi.nlm.nih.gov/myncbi/anthony.barrows.1/bibliography/public/}

	
\end{document}